\documentclass[11pt,a4paper]{article}

\usepackage[utf8]{inputenc}
\usepackage[T1]{fontenc}
\usepackage[french]{babel}
\usepackage{amsmath,amssymb,amsfonts}
\usepackage{geometry}
\usepackage{graphicx}
\usepackage{booktabs}
\usepackage{hyperref}
\usepackage{array}
\usepackage{xcolor}
\usepackage{float}

\geometry{margin=2.5cm}
\graphicspath{{figures/}}

\title{
\textbf{Courbure d'Ollivier--Ricci et limite de Ricci continue}\\
\large Vers la seconde flèche géométrique de Bottom-Up Quantum Gravity
}

\author{Farid Hamdad}
\date{\today}

\begin{document}

\maketitle

\begin{abstract}
Paper~15 a montré que la courbure d'Ollivier--Ricci discrète distingue une
géométrie positivement courbée d'une géométrie plate périodique. Le présent
papier étudie l'étape suivante : sous quelle forme la courbure
d'Ollivier--Ricci peut-elle être renormalisée pour approcher le signal de
Ricci continu
\[
R_{\mu\nu}u^\mu u^\nu ?
\]
Nous testons cette question sur deux géométries contrôlées : le tore plat
périodique, pour lequel \(R_{\mu\nu}u^\mu u^\nu=0\), et la sphère unité, pour
laquelle \(R_{\mu\nu}u^\mu u^\nu=1\). Le premier test montre que le signal
sphère--plat est positif pour la courbure brute, pour
\(\kappa^{OR}/\epsilon\), et pour \(\kappa^{OR}/\ell^2\). Le second test
identifie une calibration affine finie :
\[
\frac{\kappa^{OR}}{\epsilon}
\simeq
B_N+C_N R_{\mu\nu}u^\mu u^\nu.
\]
À \(N=512\), on obtient
\[
B_N=-0.301281,
\qquad
C_N=0.391933.
\]
Cela donne le proxy calibré
\[
\widehat{R}_{OR}
=
2.551455
\left(
\frac{\kappa^{OR}}{\epsilon}
+
0.301281
\right),
\]
qui envoie le tore plat vers une courbure moyenne \(0\) et la sphère unité
vers une courbure moyenne \(1\). Ces résultats ne prouvent pas encore une
convergence locale point par point, mais établissent un proxy moyen de Ricci
à partir de la courbure d'Ollivier--Ricci.
\end{abstract}

\tableofcontents

\section{Introduction}

Dans Bottom-Up Quantum Gravity, la géométrie émerge d'un réseau
d'intrication. L'objet fondamental est la matrice d'information mutuelle
\[
W_{ij}=I(i:j),
\]
à partir de laquelle on construit un graphe pondéré, un Laplacien
d'intrication et des observables géométriques discrètes.

Paper~16 a étudié la première flèche géométrique :
\[
L_N\longrightarrow\Delta_g.
\]
Le présent papier étudie la seconde :
\[
\kappa_{ij}^{OR}
\longrightarrow
R_{\mu\nu}u^\mu u^\nu.
\]

La courbure d'Ollivier--Ricci est une notion de courbure discrète fondée sur
le transport optimal. Elle compare la distance entre deux voisinages à la
distance entre leurs centres. Elle est donc adaptée à une géométrie
relationnelle où la notion de voisinage précède la métrique continue.

Paper~15 avait établi un premier résultat relatif :
\[
\bar\kappa_{\rm sphere}
>
\bar\kappa_{\rm flat}.
\]
Paper~17 cherche à comprendre la forme de renormalisation nécessaire pour
transformer ce signal discret en proxy de Ricci continu.

La question centrale est :
\[
\boxed{
\frac{\kappa^{OR}}{\epsilon}
=
B_N
+
C_N R_{\mu\nu}u^\mu u^\nu
+
o(1)
?
}
\]

\section{Courbure d'Ollivier--Ricci}

La courbure d'Ollivier--Ricci d'une arête \((i,j)\) est définie par :
\[
\kappa_{ij}^{OR}
=
1-
\frac{W_1(m_i,m_j)}{d(i,j)}.
\]

Ici, \(W_1(m_i,m_j)\) est la distance de Wasserstein-1 entre les mesures de
voisinage \(m_i\) et \(m_j\), et \(d(i,j)\) est la distance entre les deux
nœuds.

La mesure \(m_i\) est définie avec une inertie \(\alpha\) :
\[
m_i
=
\alpha\delta_i
+
(1-\alpha)
\sum_{k\sim i}
\frac{W_{ik}}{\sum_\ell W_{i\ell}}
\delta_k.
\]

Dans les tests numériques de ce papier, nous utilisons :
\[
\alpha=0.5.
\]

La limite continue attendue est de la forme :
\[
\kappa_{ij}^{OR}
\sim
\epsilon_N R_{\mu\nu}u^\mu u^\nu
+
\text{biais de discrétisation}.
\]

Cela motive l'étude de la quantité :
\[
\frac{\kappa_{ij}^{OR}}{\epsilon_N}.
\]

\section{Héritage de Paper 15}

Paper~15 a testé la courbure d'Ollivier--Ricci sur des géométries contrôlées.
Une première version utilisant une grille plane avec bord était contaminée par
les effets de frontière. La version corrigée a utilisé un tore plat périodique
comme référence de courbure nulle.

Les résultats finaux étaient :
\[
\bar\kappa_{\rm flat}
=
-0.001671
\simeq0,
\]
et :
\[
\bar\kappa_{\rm sphere}
=
0.008854.
\]

Le signal relatif était donc :
\[
\Delta\bar\kappa_{\rm sphere-flat}
=
0.010525.
\]

Ce résultat confirmait que la courbure d'Ollivier--Ricci distingue bien une
géométrie positivement courbée d'une géométrie plate périodique. Cependant, il
ne disait pas encore comment relier quantitativement \(\kappa^{OR}\) à
\(R_{\mu\nu}u^\mu u^\nu\).

\section{Géométries contrôlées}

Nous utilisons deux géométries de référence.

\subsection{Tore plat}

Le tore plat périodique
\[
T^2=[0,1)^2
\]
est utilisé comme référence de courbure nulle :
\[
R_{\mu\nu}u^\mu u^\nu=0.
\]

La distance intrinsèque est :
\[
d(x_i,x_j)^2
=
\min(|u_i-u_j|,1-|u_i-u_j|)^2
+
\min(|v_i-v_j|,1-|v_i-v_j|)^2.
\]

\subsection{Sphère unité}

La sphère unité \(S^2\) est utilisée comme référence de courbure positive.
Sur \(S^2\), le tenseur de Ricci vaut :
\[
{\rm Ric}=g.
\]

Ainsi, pour tout vecteur tangent unitaire \(u^\mu\),
\[
R_{\mu\nu}u^\mu u^\nu=1.
\]

La distance géodésique utilisée est :
\[
d(x_i,x_j)
=
\arccos(x_i\cdot x_j).
\]

\section{Test v1 : signal relatif sphère--plat}

Le premier test Paper~17 compare trois observables :

\[
\kappa^{OR},
\]

\[
\frac{\kappa^{OR}}{\epsilon},
\]

\[
\frac{\kappa^{OR}}{\ell^2},
\]
où \(\ell\) est la longueur intrinsèque de l'arête.

À \(N=256\), on obtient :
\[
\Delta\bar\kappa_{\rm sphere-flat}
=
0.013383.
\]

Pour la normalisation par \(\epsilon\) :
\[
\Delta
\left\langle
\frac{\kappa}{\epsilon}
\right\rangle_{\rm sphere-flat}
=
0.197719.
\]

Pour la normalisation par \(\ell^2\) :
\[
\Delta
\left\langle
\frac{\kappa}{\ell^2}
\right\rangle_{\rm sphere-flat}
=
0.557041.
\]

Les trois signaux sont positifs :
\[
\bar\kappa_{\rm sphere}
>
\bar\kappa_{\rm flat},
\]
\[
\left\langle\frac{\kappa}{\epsilon}\right\rangle_{\rm sphere}
>
\left\langle\frac{\kappa}{\epsilon}\right\rangle_{\rm flat},
\]
\[
\left\langle\frac{\kappa}{\ell^2}\right\rangle_{\rm sphere}
>
\left\langle\frac{\kappa}{\ell^2}\right\rangle_{\rm flat}.
\]

Ce résultat montre que le signal sphère--plat est robuste au choix de
normalisation. Il ne dépend pas uniquement de la courbure brute.

\section{Test v2 : calibration affine}

Le test v2 introduit une calibration affine :
\[
\widehat{R}_{OR}
=
A_N
\left(
\frac{\kappa^{OR}}{\epsilon}
-
B_N
\right).
\]

Le tore plat est utilisé pour définir le biais de courbure nulle :
\[
B_N
=
\left\langle
\frac{\kappa^{OR}}{\epsilon}
\right\rangle_{\rm flat}.
\]

La sphère unité est utilisée pour définir la normalisation :
\[
R_{\mu\nu}u^\mu u^\nu=1.
\]

On pose donc :
\[
A_N=
\frac{1}{
\left\langle\kappa^{OR}/\epsilon\right\rangle_{\rm sphere}
-
\left\langle\kappa^{OR}/\epsilon\right\rangle_{\rm flat}
}.
\]

Les constantes obtenues sont :
\[
\begin{array}{cccc}
\toprule
N_{\rm input} & B_N & \Delta_{\rm sphere-flat} & A_N\\
\midrule
128 & -0.2103 & 0.3053 & 3.2760\\
256 & -0.3079 & 0.4005 & 2.4968\\
512 & -0.3013 & 0.3919 & 2.5515\\
\bottomrule
\end{array}
\]

La stabilité entre \(N=256\) et \(N=512\) est encourageante :
\[
B_N\simeq -0.30,
\qquad
\Delta_{\rm sphere-flat}\simeq0.40,
\qquad
A_N\simeq2.5.
\]

À \(N=512\), on obtient :
\[
\widehat{R}_{OR}
=
2.551455
\left(
\frac{\kappa^{OR}}{\epsilon}
+
0.301281
\right).
\]

Par construction :
\[
\left\langle
\widehat{R}_{OR}
\right\rangle_{\rm flat}
=
0,
\]
et :
\[
\left\langle
\widehat{R}_{OR}
\right\rangle_{\rm sphere}
=
1.
\]

Cela identifie une relation moyenne :
\[
\frac{\kappa^{OR}}{\epsilon}
\simeq
B_N+C_N R_{\mu\nu}u^\mu u^\nu,
\]
avec, à \(N=512\),
\[
B_N=-0.301281,
\qquad
C_N=0.391933.
\]

\section{Résumé des résultats numériques}

Les résultats principaux sont :

\[
\begin{array}{lllll}
\toprule
{\rm Étape} & {\rm Quantité} & N & {\rm Valeur} & {\rm Statut}\\
\midrule
{\rm v1} & \Delta\bar\kappa_{\rm sphere-flat} & 256 & 0.013383 & positif\\
{\rm v1} & \Delta\langle\kappa/\epsilon\rangle_{\rm sphere-flat} & 256 & 0.197719 & positif\\
{\rm v1} & \Delta\langle\kappa/\ell^2\rangle_{\rm sphere-flat} & 256 & 0.557041 & positif\\
{\rm v2} & B_N & 512 & -0.301281 & identifié\\
{\rm v2} & C_N & 512 & 0.391933 & positif\\
{\rm v2} & A_N=1/C_N & 512 & 2.551455 & identifié\\
\bottomrule
\end{array}
\]

Le résultat central est :
\[
\boxed{
\frac{\kappa^{OR}}{\epsilon}
\simeq
-0.301
+
0.392\,R_{\mu\nu}u^\mu u^\nu
}
\]
à \(N=512\), pour les deux références géométriques utilisées.

\section{Interprétation}

Paper~17 établit que la courbure d'Ollivier--Ricci porte un signal de Ricci
relatif. La sphère présente systématiquement une courbure moyenne supérieure
au tore plat. De plus, après soustraction d'un biais de géométrie plate et
calibration par la sphère unité, la quantité \(\kappa^{OR}/\epsilon\) fournit
un proxy moyen de \(R_{\mu\nu}u^\mu u^\nu\).

Le résultat n'est pas encore une convergence locale point par point :
\[
\kappa_{ij}^{OR}
\not\equiv
\epsilon C_N R_{\mu\nu}u^\mu u^\nu
\]
pour chaque arête individuelle.

Il s'agit plutôt d'une relation moyenne :
\[
\left\langle
\frac{\kappa^{OR}}{\epsilon}
\right\rangle
\simeq
B_N+C_N
\left\langle
R_{\mu\nu}u^\mu u^\nu
\right\rangle.
\]

C'est néanmoins une étape importante, car elle transforme le signal qualitatif
de Paper~15 en relation quantitative calibrée.

\section{Limites}

Plusieurs limites doivent être précisées.

Premièrement, la calibration affine est empirique. Les constantes \(B_N\) et
\(C_N\) ne sont pas encore dérivées analytiquement.

Deuxièmement, le résultat est un proxy moyen. La dispersion des valeurs de
\(\widehat{R}_{OR}\) reste importante, notamment sur le tore plat.

Troisièmement, le graphe \(k\)-NN peut introduire une anisotropie locale. Une
version ultérieure devra tester des graphes à rayon fixe.

Quatrièmement, les tests ne concernent que deux géométries :
\[
T^2,
\qquad
S^2.
\]
Il faudra tester des sphères de rayon différent, des surfaces hyperboliques et
des géométries mixtes.

Cinquièmement, la relation avec de vrais graphes d'information mutuelle
\[
W_{ij}=I(i:j)
\]
reste à établir. Les tests actuels utilisent des noyaux géométriques
contrôlés.

\section{Lien avec l'équation d'Einstein effective}

Paper~15 a proposé la chaîne :
\[
W_{ij}
\to
L_\epsilon
\to
\Delta_g,
\]
\[
\kappa_{ij}^{OR}
\to
R_{\mu\nu}u^\mu u^\nu,
\]
\[
\delta W_{\rm loc}
\to
\delta\kappa(r).
\]

Paper~16 a renforcé la première flèche au niveau spectral. Paper~17 renforce
la seconde en montrant qu'une calibration affine de \(\kappa^{OR}/\epsilon\)
peut reproduire les valeurs moyennes attendues de Ricci sur deux géométries
contrôlées.

Ainsi, dans le programme complet :
\[
\frac{\delta S_{\rm BuP}}{\delta W_{ij}}=0
\quad
\Longrightarrow
\quad
G_{\mu\nu}
+
\Lambda_{\rm ent}g_{\mu\nu}
=
8\pi G_{\rm eff}T_{\mu\nu}^{\rm ent}
+
\mathcal{H}_{\mu\nu},
\]
Paper~17 fournit une brique pour identifier le contenu continu du terme de
courbure.

\section{Conclusion}

Paper~17 montre que la courbure d'Ollivier--Ricci discrète contient un signal
de Ricci exploitable. Le signal sphère--plat est positif pour la courbure
brute, pour \(\kappa/\epsilon\), et pour \(\kappa/\ell^2\). La calibration
affine v2 identifie ensuite une relation moyenne :
\[
\frac{\kappa^{OR}}{\epsilon}
\simeq
B_N+C_N R_{\mu\nu}u^\mu u^\nu.
\]

À \(N=512\), cette relation prend la forme :
\[
\frac{\kappa^{OR}}{\epsilon}
\simeq
-0.301
+
0.392\,R_{\mu\nu}u^\mu u^\nu.
\]

Le proxy calibré est :
\[
\widehat{R}_{OR}
=
2.551
\left(
\frac{\kappa^{OR}}{\epsilon}
+
0.301
\right).
\]

Il envoie le tore plat vers une courbure moyenne \(0\) et la sphère unité vers
une courbure moyenne \(1\). Ce résultat ne prouve pas encore la convergence
locale :
\[
\kappa_{ij}^{OR}
\to
R_{\mu\nu}u^\mu u^\nu,
\]
mais il fournit la première calibration quantitative du signal de Ricci dans
BuP.

\section*{Disponibilité du code et des données}

Les scripts et résultats de Paper~17 sont organisés dans le dossier :
\[
\texttt{papers/paper17\_ollivier\_ricci\_limit/}.
\]

Les scripts principaux sont :
\[
\texttt{scripts/paper17\_or\_ricci\_scaling\_v1.py},
\]
\[
\texttt{scripts/paper17\_or\_ricci\_scaling\_v2.py},
\]
\[
\texttt{scripts/paper17\_build\_ricci\_summary\_v1.py}.
\]

Les résultats principaux sont stockés dans :
\[
\texttt{results/or\_ricci\_scaling\_v1/},
\]
\[
\texttt{results/or\_ricci\_scaling\_v2/},
\]
\[
\texttt{results/paper17\_ricci\_summary\_v1/}.
\]

\appendix

\section{Détails sur la calibration}

La calibration affine est définie par :
\[
\widehat{R}_{OR}
=
A_N
\left(
\frac{\kappa^{OR}}{\epsilon}
-
B_N
\right).
\]

Le biais plat est :
\[
B_N=
\left\langle
\frac{\kappa^{OR}}{\epsilon}
\right\rangle_{\rm flat}.
\]

La normalisation est :
\[
A_N=
\frac{1}{
\left\langle\kappa^{OR}/\epsilon\right\rangle_{\rm sphere}
-
\left\langle\kappa^{OR}/\epsilon\right\rangle_{\rm flat}
}.
\]

En posant :
\[
C_N=\frac{1}{A_N},
\]
on obtient :
\[
\frac{\kappa^{OR}}{\epsilon}
\simeq
B_N+C_N R_{\mu\nu}u^\mu u^\nu.
\]

\section{Prochaines étapes}

Les prochaines étapes sont :

\begin{enumerate}
    \item tester des graphes à rayon fixe au lieu de \(k\)-NN ;
    \item réduire la dispersion du tore plat ;
    \item binning par longueur d'arête \(\ell_{ij}\) ;
    \item utiliser des moyennes tronquées ;
    \item augmenter \(N\) avec transport optimal optimisé ;
    \item tester des sphères de rayon \(r\), avec
    \[
    R_{\mu\nu}u^\mu u^\nu=\frac{1}{r^2};
    \]
    \item tester des surfaces à courbure négative ;
    \item dériver analytiquement les constantes \(B_N\) et \(C_N\).
\end{enumerate}

\end{document}
